\subsection{Tzanio, MFEM}

The MFEM team will contribute a targeted ICF miniapp designed primarily as a high-fidelity data generator for training and validating the proposed surrogate modeling framework. Building on existing capabilities in Lagrangian hydrodynamics \cite{Dobrev2012}, ALE methods \cite{Tomov2018}, and scalable H(div) radiation diffusion solvers \cite{Pazner2025}, this effort will focus on producing physically consistent, multi-physics datasets that capture coupled hydrodynamics–radiation behavior across parameterized geometries and operating conditions. The miniapp will be tightly integrated into the AI workflow to systematically generate training data, support active learning loops, and enable evaluation of surrogate accuracy and generalization. In Phase I, simplified ICF configurations will be used to demonstrate feasibility and establish a robust pipeline connecting MFEM-generated data with the foundation models, ensuring that learned representations remain consistent with the underlying physical operators and constraints.

This effort directly complements the proposed AI architecture by providing high-order, structure-preserving finite element data that aligns with the project’s emphasis on physically consistent latent spaces. MFEM’s support for complex, unstructured meshes and high-order discretizations enables systematic variation of geometry and resolution, which is critical for training geometry-aware models such as GNN-based encoders. Moreover, the compatibility of MFEM operators with differentiable programming frameworks enables a natural pathway to incorporate operator-level sensitivities into both surrogate training and downstream optimization. This creates a strong link between the simulation stack and the latent-space optimization strategies outlined in the proposal, particularly for adjoint-informed refinement and hybrid optimization workflows.

In Phase II, the MFEM miniapp will be extended to more realistic ICF scenarios, including fully three-dimensional configurations, advanced material models, and increasingly complex mesh topologies. These enhancements will significantly enrich the training corpus, improving the robustness and generalization of the surrogate models across regimes relevant to fusion energy applications. The MFEM GPU capabilities \cite{MFEM2024} will be leveraged to scale data generation efficiently on next-generation DOE platforms, aligning with the Genesis infrastructure. In addition, recently developed differentiability capabilities within MFEM will be used to expose gradients of the finite element operators, enabling tighter integration with optimization and control components of the framework. This will support gradient-informed training, accelerate convergence of surrogate models, and provide a principled mechanism to couple high-fidelity physics with the agentic optimization strategies envisioned for Phase II.
